% !TeX root = ../main.tex
% ---------------------------------------------------------------------------
% Multiple Excavations, Pillars and Stress Interaction
% Scope extension added from the classical-textbook crosscheck.
% Maps to: Hoek & Brown (1980) Ch. 5 (multiple excavations, pillars, stress
% shadows); Brady & Brown (2006) Ch. 13 "Pillar supported mining
% methods" and Ch. 10 (pillar crushing).
% Rationale: the book treats single openings throughout. Twin tunnels,
% station caverns, cross-passages and multi-level works are all interaction
% problems, and no chapter currently covers pillar behaviour.
% Mining-specific content is deliberately adapted, not reproduced.
% ---------------------------------------------------------------------------
\chapter{Multiple Excavations, Pillars and Stress Interaction}
\label{ch:pillars-interaction}


\section{Interaction between openings}
\label{sec:interaction-between-openings}

\subsection{Openings rarely occur alone}
\label{sec:openings-interact}
\scopenote{Twin running tunnels, station caverns with adjacent concourses,
cross-passages, ventilation adits, stacked alignments and works adjacent to
existing tunnels. Interaction is the normal case, not the exception.}

\subsection{Superposition of stress fields}
\label{sec:stress-superposition}
\scopenote{When elastic superposition is admissible, and why it fails once
yielding zones from adjacent openings overlap.}

\subsection{Zone of influence and minimum separation}
\label{sec:separation-criteria}
\scopenote{Using the zone of influence from
Section~\ref{sec:zone-of-influence} to set separation rules of thumb, and the
conditions under which the rules are unsafe.}

\section{Stress shadows and protected zones}
\label{sec:stress-shadows}
\scopenote{The de-stressed region created by an adjacent opening, its
exploitation in sequencing, and the loss of confinement that accompanies it.}

\section{Pillar stress}
\label{sec:pillar-stress}

\subsection{Pillar stress by tributary area}
\label{sec:tributary-area}
\scopenote{The simplest pillar stress estimate, its assumptions of uniform
extraction and rigid abutments, and the extraction ratio relationship.}

\subsection{Numerical determination of pillar stress}
\label{sec:numerical-pillar-stress}
\scopenote{Where tributary area breaks down --- irregular layouts, varying
depth, few pillars --- and how numerical analysis resolves the actual stress
distribution across a pillar.}

\section{Pillar strength and design}
\label{sec:pillar-strength-and-design}

\subsection{Pillar strength}
\label{sec:pillar-strength}
\scopenote{Empirical strength formulae, the width-to-height ratio effect,
confinement in the pillar core, and the size effect that makes a large pillar
weaker per unit area than a laboratory specimen.}

\subsection{Pillar factor of safety and design}
\label{sec:pillar-fos}
\scopenote{Deterministic and probabilistic factor of safety, typical target
values by pillar duty, and design charts.}

\section{Pillar failure}
\label{sec:pillar-failure-group}

\subsection{Pillar failure modes}
\label{sec:pillar-failure}
\scopenote{Spalling and progressive slabbing, shear through weak planes,
foundation failure and squeezing; the distinction between a yielding pillar
that still carries load and one that has failed.}

\subsection{Progressive and domino failure}
\label{sec:progressive-failure}
\scopenote{Load transfer on pillar failure, the conditions for cascading
collapse, and the stiffness relationship between pillar and surrounding rock
mass that determines whether failure is stable or violent.}

\subsection{Roof and floor bearing capacity}
\label{sec:pillar-foundations}
\scopenote{Punching of a strong pillar into a weak roof or floor, and the
often-overlooked case where the pillar is not the weakest element.}

\section{Applications in civil works}
\label{sec:applications-in-civil-works}

\subsection{Twin tunnel pillars in civil works}
\label{sec:twin-tunnel-pillars}
\scopenote{The rock or soil pillar between running tunnels: minimum pillar
width, the influence of construction sequence, and the settlement consequences
covered in Chapter~\ref{ch:settlement-damage}.}

\subsection{Multi-cavern and station layouts}
\label{sec:multi-cavern}
\scopenote{Powerhouse and transformer hall pairs, station cavern and adit
clusters, and layout optimisation against both stress interaction and
functional requirements.}

\subsection{Excavation sequence and interaction}
\label{sec:sequence-interaction}
\scopenote{How the order of excavation changes the stress path experienced by
the pillar and the final support demand, and why sequence must be specified in
the design, not left to the contractor.}
